%------------------------------------------------------------------------------%
\begin{exercises*}

% Exercício de prova do prof. André Pereira
\begin{exercise}
  Considere a série de potências
  \[
    \sum_{n=0}^\infty \frac{(-1)^n(5x-5)^{n}}{n5^{n}}
  \]
  \vspace{-\baselineskip}
  \begin{alphalist}
    \item Qual o centro $a$ da série?
    \item Qual o raio de convergência?
    \item Qual o intervalo de convergência da série?
    \item Encontre uma função $f(x)$ que coincide com a série no interior do intervalo de convergência.
  \end{alphalist}
\end{exercise}

% Exercício de prova do prof. André Pereira
\begin{exercise}
  Considere a série de potências
  \[\sum_{n=1}^\infty (-1)^{n-1}\frac{x^{2n-1}}{2n-1}\]
  \begin{alphalist}
    \item Encontre o intervalo de convergência da série.
    \item Encontre uma função $f$ tal que
      \[ f(x)=\sum_{n=0}^\infty (-1)^{n-1}\frac{x^{2n-1}}{2n-1}\]
    no interior do intervalo de convergência.
  \end{alphalist}
\end{exercise}

%    \begin{exercise}
%        \begin{mathlist}[2]
%            \item
%            \item
%            \item
%        \end{mathlist}
%        \begin{answer}
%            \begin{alphalist}[3]
%                \item
%                \item
%                \item
%            \end{alphalist}
%        \end{answer}
%    \end{exercise}

\end{exercises*}
%------------------------------------------------------------------------------%
